\subsection{RQ5: reference points}
\label{sec:rq_reference}

Table~\ref{tab:baselines} places our views next to re-implementations of the
\emph{input representations} used by existing approaches, a pretrained code
model, and the trivial size rule, all on the same samples and folds.
They are not reproductions of the published systems or their published numbers,
which use different data splits and often within-project evaluation.

% Generated by experiments/emit_tables.py -- DO NOT EDIT BY HAND.
% Every value is read from experiments/results/*.json.

\begin{table*}[t]
\centering
\caption{Comparison with existing approaches re-implemented under the same
sample, folds and metrics (mean $\pm$ std over 10 project-level folds). These
are re-implementations of the \emph{input representation} of each approach
inside one training protocol, not reproductions of the published systems or
their published numbers. UniXcoder rows use a frozen encoder and a linear
probe. \emph{Format-normalised} source is the parser's token stream with
whitespace and comments removed (Section~\ref{sec:artefact}). The last three
rows score functions without reading what the code says.}
\label{tab:baselines}
\resizebox{\linewidth}{!}{%
\begin{tabular}{lcccc}
\toprule
 & \multicolumn{2}{c}{BigVul} & \multicolumn{2}{c}{DiverseVul} \\
System & AUC & AUPRC & AUC & AUPRC \\
\midrule
GGNN on code text (token-only baseline) & 0.506 $\pm$ 0.077 & 0.085 $\pm$ 0.034 & 0.566 $\pm$ 0.039 & 0.080 $\pm$ 0.017 \\
Devign/ReVeal-style GGNN: text $+$ node type & 0.592 $\pm$ 0.072 & 0.105 $\pm$ 0.039 & 0.607 $\pm$ 0.027 & 0.096 $\pm$ 0.024 \\
GGNN, anonymised tokens $+$ node type & 0.701 $\pm$ 0.059 & 0.180 $\pm$ 0.059 & 0.679 $\pm$ 0.036 & 0.122 $\pm$ 0.030 \\
Transformer over node sequence, text & 0.582 $\pm$ 0.074 & 0.106 $\pm$ 0.042 & --- & --- \\
Transformer over node sequence, $\kappa$ & 0.706 $\pm$ 0.039 & 0.155 $\pm$ 0.041 & --- & --- \\
Pretrained UniXcoder, frozen, linear probe (raw source) & 0.705 $\pm$ 0.061 & 0.163 $\pm$ 0.035 & 0.623 $\pm$ 0.027 & 0.095 $\pm$ 0.022 \\
Pretrained UniXcoder, frozen, linear probe (format-normalised) & 0.645 $\pm$ 0.069 & 0.109 $\pm$ 0.022 & 0.604 $\pm$ 0.015 & 0.088 $\pm$ 0.019 \\
LineVul-style: fine-tuned CodeBERT (raw source) & 0.922 $\pm$ 0.026 & 0.748 $\pm$ 0.063 & 0.693 $\pm$ 0.028 & 0.123 $\pm$ 0.026 \\
LineVul-style: fine-tuned CodeBERT (format-normalised) & 0.728 $\pm$ 0.058 & 0.189 $\pm$ 0.058 & 0.705 $\pm$ 0.028 & 0.131 $\pm$ 0.023 \\
No code read: function size (node count) & 0.736 $\pm$ 0.065 & 0.232 $\pm$ 0.087 & 0.715 $\pm$ 0.028 & 0.152 $\pm$ 0.028 \\
No code read: trailing-whitespace rule & 0.830 $\pm$ 0.050 & 0.392 $\pm$ 0.109 & 0.542 $\pm$ 0.024 & 0.069 $\pm$ 0.018 \\
No code read: gradient boosting on formatting statistics & 0.896 $\pm$ 0.026 & 0.601 $\pm$ 0.065 & 0.634 $\pm$ 0.033 & 0.111 $\pm$ 0.034 \\
GGNN, grammar kind $\kappa$ (ours) & 0.691 $\pm$ 0.037 & 0.149 $\pm$ 0.027 & 0.667 $\pm$ 0.021 & 0.105 $\pm$ 0.026 \\
GGNN, $\kappa{+}\phi$ (ours) & 0.687 $\pm$ 0.060 & 0.158 $\pm$ 0.042 & 0.675 $\pm$ 0.021 & 0.114 $\pm$ 0.030 \\
\bottomrule
\end{tabular}}
\end{table*}


\paragraph{Devign/ReVeal-style input} Text together with the node type
(\texttt{lex+$\kappa$} with a GGNN) reaches AUC \mBvGgnnLexKindAuc{} on BigVul
and \mDvGgnnLexKindAuc{} on DiverseVul, which is below the grammar-kind and anonymised views in
Table~\ref{tab:rq2_representation}.

\paragraph{A pretrained code model} UniXcoder~\cite{guo2022unixcoder} with a frozen
encoder (its weights are not updated; it turns each function into one vector) and
a linear probe (a logistic regression trained on that vector) on the raw source (identifiers included) reaches AUC
\numbvUnixAuc{} on BigVul and \numdvUnixAuc{} on DiverseVul. On BigVul it is
statistically indistinguishable from the vocabulary-free graph views (the paired
difference of $\kappa$ against it is \cBvGgnnKindVUnixD{} with the GGNN,
\cBvGgnnKindVUnixLo{} to \cBvGgnnKindVUnixHi{}, and similar for the other
backbones); on DiverseVul the grammar-kind GNNs are better (GGNN
\cDvGgnnKindVUnixD{}, \cDvGgnnKindVUnixLo{} to \cDvGgnnKindVUnixHi{}; GIN
\cDvGinKindVUnixD{}; GAT \cDvGatKindVUnixD{}, not significant). Two lessons
follow. First, ``text hurts cross-project transfer'' is a statement about token
embeddings learned from scratch on a few thousand functions, not about text in
general: a model that has seen billions of code tokens before can use
identifiers and still transfer about as well. Second, a vocabulary-free GNN
trained from scratch is comparable with, not clearly better than, a frozen
pretrained model. The probe is frozen and untuned; a fine-tuned LineVul-style
model~\cite{fu2022linevul} would probably be stronger, and we have not run one
because this machine has no GPU.

\paragraph{A size rule has the highest mean AUC} Scoring a function by its node
count alone gives AUC \numbvSizeAuc{} (BigVul) and \numdvSizeAuc{} (DiverseVul),
the highest mean AUC in this study. The paired differences are significant under
the project-cluster bootstrap for most networks (e.g.\ $\kappa$ with the GGNN
against size: \cBvGgnnKindVSizeD{}, \cBvGgnnKindVSizeLo{} to
\cBvGgnnKindVSizeHi{} on BigVul) but not for all (for example the GIN with
$\kappa{+}\phi$ on BigVul). We read this as a caution about the whole
benchmark family: under project-disjoint evaluation, much of what these
corpora reward is a length prior, and absolute AUCs of $0.6$ to $0.7$ should not
be read as evidence of vulnerability understanding.
